\unit{Exposição 195 -- Técnica de descida II}
\sid{195}
\src{369}
\seminar{12.\textsuperscript{o} ano, 1959/60, n.\textsuperscript{o} 195\\Fevereiro de 1960}
\paperhead{TÉCNICA DE DESCIDA E TEOREMAS DE EXISTÊNCIA EM GEOMETRIA ALGÉBRICA\\
II. O TEOREMA DE EXISTÊNCIA NA TEORIA FORMAL DOS MÓDULOS}%
  {por Alexander GROTHENDIECK}

\fsection{A}{Funtores representáveis e pro-representáveis}

\subsection*{1. Funtores representáveis}
\addcontentsline{toc}{subsection}{1. Funtores representáveis}

Seja $\mathcal C$ uma categoria. Para todo $X\in\mathcal C$, seja $h_X$ o
funtor contravariante de $\mathcal C$ na categoria $(\operatorname{Ens})$
dos conjuntos,
\[
  h_X:\mathcal C\longrightarrow(\operatorname{Ens}),
\]
definido pela fórmula
\[
  h_X(Y)=\operatorname{Hom}(Y,X).
\]
Um morfismo $X\longrightarrow X'$ em $\mathcal C$ induz de maneira evidente um
homomorfismo de funtores $h_X\longrightarrow h_{X'}$: $h_X$ é um funtor
covariante em $X$; isto é, temos definido um funtor covariante canônico
\[
  h:\mathcal C\longrightarrow
  \operatorname{Hom}\bigl(\mathcal C^o,(\operatorname{Ens})\bigr)
\]
de $\mathcal C$ na categoria dos funtores covariantes da categoria dual
$\mathcal C^o$ de $\mathcal C$ na categoria dos conjuntos. Recordemos
então:

\result{PROPOSIÇÃO}{1.1} Este funtor $h$ é \emph{plenamente fiel}; em outros
termos, para todo par $X,X'$ de objetos de $\mathcal C$, a aplicação
natural
\[
  \operatorname{Hom}(X,X')\longrightarrow
  \operatorname{Hom}(h_X,h_{X'})
\]
é bijetiva.

Em particular, se um funtor
$F\in\operatorname{Hom}(\mathcal C^o,(\operatorname{Ens}))$ é isomorfo a um
funtor da forma $h_X$, então $X$ fica determinado salvo isomorfismo
único. Diz-se nesse caso que o funtor $F$ é \emph{representável}. A
proposição anterior significa que o funtor canônico $h$ define uma
equivalência da categoria $\mathcal C$ com a subcategoria plena de
$\operatorname{Hom}(\mathcal C^o,(\operatorname{Ens}))$ formada pelos
funtores representáveis. Este fato fundamenta a noção de «solução de um
problema universal»: tal problema consiste sempre em examinar se um funtor
dado (contravariante como aqui, ou covariante no caso dual) de $\mathcal C$
em $(\operatorname{Ens})$ é representável.

\src{370}
Observemos além disso que, pela própria definição de produto em uma categoria
[1], o funtor $h:X\longmapsto h_X$ comuta com os produtos sempre que estes
existam (e, mais geralmente, com os limites projetivos finitos ou infinitos;
em particular, com os produtos fibrados, com a formação de «núcleos» [2],
etc., sempre que existam): tem-se um isomorfismo de funtores
\[
  h_{X\times X'}\xrightarrow{\sim}h_X\times h_{X'}
\]
sempre que exista $X\times X'$; isto é, têm-se bijeções funtoriais
em $Y$
\[
  h_{X\times X'}(Y)\xrightarrow{\sim}h_X(Y)\times h_{X'}(Y).
\]
Em particular, dar um morfismo
\[
  X\times X'\longrightarrow X''
\]
em $\mathcal C$ (isto é, uma «lei de composição» em $\mathcal C$ entre
$X,X',X''$) equivale a dar um morfismo
\[
  h_{X\times X'}=h_X\times h_{X'}\longrightarrow h_{X''},
\]
isto é, a dar, para todo $Y\in\mathcal C$, uma lei de composição de
conjuntos
\[
  h_X(Y)\times h_{X'}(Y)\longrightarrow h_{X''}(Y)
\]
de tal modo que, para todo morfismo $Y\longrightarrow Y'$ em $\mathcal C$, o
sistema de aplicações entre conjuntos
\[
  h_{X^{(i)}}(Y)\longleftarrow h_{X^{(i)}}(Y')
  \qquad(\text{para }i=0,1,2)
\]
seja um morfismo com respeito às duas leis de composição correspondentes a
$Y$ e $Y'$. Vê-se assim que a noção de estrutura de «$\mathcal C$-grupo»,
«$\mathcal C$-anel», etc. sobre um objeto $X$ de $\mathcal C$ se exprime do
modo mais cômodo (tanto na teoria como na prática) dizendo que, para
todo $Y\in\mathcal C$, o conjunto $h_X(Y)$ está provido de uma lei de grupo
(resp. de anel, etc.) no sentido usual, e que as aplicações
$h_X(Y)\longleftarrow h_X(Y')$ correspondentes a morfismos
$Y\longrightarrow Y'$ são homomorfismos de grupos (resp. de anéis, etc.).
Esta é, por exemplo, a forma mais intuitiva e cômoda de definir os diversos
grupos clássicos $G_a$, $G_m$, $\operatorname{Gl}(n)$, etc. sobre um pré-esquema
base arbitrário $S$ e de escrever as relações clássicas entre esses grupos;
de dotar o esquema afim $V(\mathcal F)$ sobre $S$, definido por um feixe
quase coerente $\mathcal F$, de uma estrutura de «fibrado vetorial»; e de
definir e estudar as diversas variedades de bandeiras associadas
(grassmannianas, fibrados projetivos), etc. O yoga geral consiste
simplesmente em identificar, mediante o funtor canônico $h$, os objetos de
$\mathcal C$ com certos funtores contravariantes: os funtores representáveis,

\src{371}
de $\mathcal C$ na categoria dos conjuntos.

O procedimento habitual de inverter as setas, necessário, por exemplo, no
caso dos esquemas afins para passar da linguagem geométrica à da álgebra
comutativa, leva a dualizar as considerações anteriores e, em
particular, a introduzir também os funtores covariantes representáveis
$\mathcal C\longrightarrow(\operatorname{Ens})$; isto é, os da forma
\[
  Y\longmapsto\operatorname{Hom}(X,Y)=h'_X(Y).
\]

\subsection*{2. Funtores pro-representáveis, pro-objetos}
\addcontentsline{toc}{subsection}{2. Funtores pro-representáveis, pro-objetos}

Seja $\widetilde X=(X_i)_{i\in I}$ um sistema projetivo de objetos de
$\mathcal C$; a ele corresponde um funtor covariante
\[
  h'_{\widetilde X}=\varinjlim_i h'_{X_i},
\]
mais explicitamente,
\[
  h'_{\widetilde X}(Y)=\varinjlim_i h'_{X_i}(Y)
  =\varinjlim_i\operatorname{Hom}(X_i,Y),
\]
de $\mathcal C$ em $(\operatorname{Ens})$. Um funtor de $\mathcal C$ em
$(\operatorname{Ens})$ que seja isomorfo a um funtor desse tipo com $I$
filtrante chama-se \emph{pro-representável}. Pelo número anterior, estes são
exatamente os funtores isomorfos a limites indutivos filtrantes de funtores
representáveis. Seja $\widetilde X'=(X'_j)_{j\in J}$ um segundo sistema
projetivo filtrante em $\mathcal C$ (construído sobre um segundo conjunto
preordenado filtrante de índices, $J$); verifica-se facilmente que há uma
bijeção canônica (que generaliza a proposição 1.1)
\[
  \operatorname{Hom}(h'_{\widetilde X'},h'_{\widetilde X})
  =\varprojlim_j\varinjlim_i\operatorname{Hom}(X_i,X'_j).
\]
Isso leva a introduzir a categoria $\operatorname{Pro}(\mathcal C)$ dos
pro-objetos de $\mathcal C$; seus objetos são os sistemas projetivos de objetos
de $\mathcal C$ (sobre conjuntos variáveis de índices preordenados filtrantes),
e se $\widetilde X=(X_i)_{i\in I}$ e
$\widetilde X'=(X'_j)_{j\in J}$ são dois objetos dessa classe, põe-se
\[
  \operatorname{ProHom}(\widetilde X,\widetilde X')
  =\varprojlim_j\varinjlim_i\operatorname{Hom}(X_i,X'_j),
\]
sendo evidente a composição dos pro-homomorfismos. Por sua própria
construção, $\widetilde X\longmapsto h'_{\widetilde X}$ pode ser considerado
como um funtor contravariante em $\widetilde X$, que estabelece uma equivalência
entre a categoria dual da categoria $\operatorname{Pro}(\mathcal C)$ dos
pro-objetos de $\mathcal C$ e a categoria dos funtores covariantes
pro-representáveis de $\mathcal C$ em $(\operatorname{Ens})$. Naturalmente, um
objeto $X$ de $\mathcal C$ define canonicamente um pro-objeto, denotado

\src{372}
também por $X$, de modo que $\mathcal C$ é equivalente a uma subcategoria
plena de $\operatorname{Pro}(\mathcal C)$. Se então
$\widetilde X=(X_i)_{i\in I}$ é um pro-objeto qualquer de $\mathcal C$, tem-se
(com a identificação anterior)
\[
  \widetilde X=\varprojlim_i X_i,
\]
tomando-se o limite projetivo em $\operatorname{Pro}(\mathcal C)$
(visto que $h'_{\widetilde X}=\varinjlim_i h'_{X_i}$). Convém observar que,
mesmo quando o limite projetivo dos $X_i$ exista em $\mathcal C$, em
geral não será isomorfo ao limite projetivo $\widetilde X$ em
$\operatorname{Pro}(\mathcal C)$, como já resulta evidente quando $\mathcal C$
é a categoria dos conjuntos. Note-se que, por definição,
$\varprojlim_{\mathcal C}X_i=L$ fica definido pela condição de que o
funtor em $Y\in\mathcal C$
\[
  \varprojlim_i\operatorname{Hom}_{\mathcal C}(Y,X_i)
  =\operatorname{Hom}_{\operatorname{Pro}(\mathcal C)}(Y,\widetilde X)
\]
com valores em $(\operatorname{Ens})$ seja representável mediante $L$, isto é,
seja isomorfo a $\operatorname{Hom}_{\mathcal C}(Y,L)$; por conseguinte,
$\varprojlim_{\mathcal C}X_i$ já está definido em termos do pro-objeto
$\widetilde X$ e, mais precisamente, depende funtorialmente do pro-objeto
$\widetilde X$ sempre que esteja definido; não há, pois, inconveniente em
denotá-lo por $\varprojlim_{\mathcal C}(\widetilde X)$. Se os limites
projetivos existem sempre em $\mathcal C$,
$\varprojlim_{\mathcal C}(\widetilde X)$ é um funtor de
$\operatorname{Pro}(\mathcal C)$ em $\mathcal C$, e há um homomorfismo
canônico de funtores
$\varprojlim_{\mathcal C}(\widetilde X)\longrightarrow\widetilde X$.

Como todo funtor, que suporemos covariante para fixar as ideias,
\[
  F:\mathcal C\longrightarrow\mathcal C'
\]
se prolonga de maneira evidente a um funtor
\[
  \operatorname{Pro}(F):\operatorname{Pro}(\mathcal C)
  \longrightarrow\operatorname{Pro}(\mathcal C'),
\]
segue-se que, se em $\mathcal C'$ sempre existem os limites projetivos, $F$
define também canonicamente um funtor composto
\[
  \overline F:\operatorname{Pro}(\mathcal C)\longrightarrow\mathcal C',
  \qquad
  \overline F=\varprojlim_{\mathcal C'}\circ\operatorname{Pro}(F),
\]
que transforma $\widetilde X=(X_i)_{i\in I}$ em
$\varprojlim_{\mathcal C'}F(X_i)$.

Um pro-objeto $\widetilde X$ chama-se \emph{pro-objeto estrito} se é isomorfo a
um pro-objeto $(X_i)_{i\in I}$ cujos morfismos de transição
$X_i\longrightarrow X_j$ são epimorfismos; um funtor definido por tal objeto
chama-se \emph{estritamente pro-representável}. Pode-se exigir, além disso, que $I$
seja ordenado filtrante e que todo epimorfismo $X_i\longrightarrow X'$ seja
equivalente a um epimorfismo

\src{373}
$X_i\longrightarrow X_j$ para um $j\in I$ adequado (determinado de maneira
única por essa condição). Nessas condições, o sistema projetivo
$(X_i)_{i\in I}$ fica determinado salvo isomorfismo único (no sentido usual
dos isomorfismos de sistemas projetivos). Segue-se que em
$\operatorname{Pro}(\mathcal C)$ existe sempre o limite projetivo de um
sistema projetivo de pro-objetos estritos $\widetilde X^{(\alpha)}$ e que, com
as notações anteriores para $F$, $\overline F$, ter-se-á
\[
  \overline F\left(\varprojlim_\alpha\widetilde X^{(\alpha)}\right)
  =\varprojlim_{\alpha,\mathcal C'}
  \overline F\left(\widetilde X^{(\alpha)}\right).
\]
Em particular, se todo pro-objeto de $\mathcal C$ é estrito (cf. o número
seguinte), o funtor prolongado $\overline F$ comuta com os limites
projetivos.

\subsection*{3. Caracterização dos funtores pro-representáveis}
\addcontentsline{toc}{subsection}{3. Caracterização dos funtores pro-representáveis}

Sejam $\mathcal C$, $\mathcal C'$ duas categorias nas quais existem os limites
projetivos finitos (isto é, relativos a conjuntos preordenados finitos, não
necessariamente filtrantes), ou, o que é equivalente, nas quais existem os
produtos finitos e os produtos fibrados finitos (o que implica, em
particular, a existência de um «objeto final» $e$, tal que
$\operatorname{Hom}(X,e)$ se reduza a um elemento para todo $X$). Seja $F$ um
funtor covariante de $\mathcal C$ em $\mathcal C'$. As condições seguintes
são equivalentes:

\begin{enumerate}
\item[(i)] $F$ comuta com os limites projetivos finitos;
\item[(ii)] $F$ comuta com os produtos finitos e os produtos fibrados finitos;
\item[(iii)] $F$ comuta com os produtos finitos e, para todo diagrama exato
\[
  X\longrightarrow X'\rightrightarrows X''
\]
em $\mathcal C$ ([3], A, definição 2.1), o diagrama transformado por $F$
\[
  F(X)\longrightarrow F(X')\rightrightarrows F(X'')
\]
é exato.
\end{enumerate}

Diz-se então que $F$ é \emph{exato à esquerda}.

De agora em diante, supõe-se que em $\mathcal C$ existem os limites projetivos
finitos. Das definições resulta então imediatamente que um funtor
representável é exato à esquerda e, passando ao limite, que um funtor
pro-representável é exato à esquerda.

Para obter uma recíproca, seja
\[
  F:\mathcal C\longrightarrow(\operatorname{Ens})
\]
um funtor covariante, e sejam $X\in\mathcal C$ e $\xi\in F(X)$. Diz-se que
$\xi$ (ou o par \src{374}$(X,\xi)$) é \emph{minimal} se, para todo
$X'\in\mathcal C$ e $\xi'\in F(X')$, e todo monomorfismo estrito ([3], A, 2)
$u:X'\longrightarrow X$ tal que $\xi=F(u)(\xi')$, $u$ é um isomorfismo. Por
outro lado, diz-se que um par $(X,\xi)$ \emph{domina} $(X'',\xi'')$
($\xi\in F(X)$, $\xi''\in F(X'')$) se existe um morfismo
$v:X\longrightarrow X''$ tal que $\xi''=F(v)(\xi)$; se $\xi$ é minimal e $F$
é exato à esquerda, este morfismo $v$ é então único; se $\xi''$ é
minimal, $v$ é sobrejetivo. Deduz-se facilmente a proposição seguinte:

\result{PROPOSIÇÃO}{3.1} Para que $F$ seja estritamente pro-representável, é
necessário e suficiente que satisfaça as duas condições:
\begin{enumerate}
\item[(i)] $F$ é exato à esquerda;
\item[(ii)] todo par $(X,\xi)$, com $\xi\in F(X)$, é dominado por um par
minimal.
\end{enumerate}

Esta última condição é vazia se todo objeto de $\mathcal C$ é artiniano
(tomando um subobjeto $X'$ de $X$ minimal entre aqueles para os quais existe
$\xi'\in F(X')$ tal que $\xi$ seja a imagem de $\xi'$). Daí:

\result{COROLÁRIO}{} Seja $\mathcal C$ uma categoria cujos objetos são
artinianos e na qual existem os limites projetivos finitos. Então os
funtores pro-representáveis de $\mathcal C$ em $(\operatorname{Ens})$ são
exatamente os funtores exatos à esquerda, e são inclusive estritamente
pro-representáveis.

Este último fato significa também que \emph{todo pro-objeto de $\mathcal C$
é estrito}.

\subsection*{4. Exemplo: grupos de tipo galoisiano, grupos proalgébricos}
\addcontentsline{toc}{subsection}{4. Grupos de tipo galoisiano, grupos proalgébricos}

Se $\mathcal C$ é a categoria dos grupos finitos ordinários,
$\operatorname{Pro}(\mathcal C)$ é equivalente à categoria dos grupos
topológicos compactos totalmente desconexos. São grupos desse tipo e suas
generalizações, obtidas substituindo os grupos finitos ordinários por
esquemas em grupos finitos sobre um pré-esquema base dado (por exemplo, os
grupos algébricos finitos sobre um corpo $k$), que farão as vezes de
grupos fundamentais, de homotopia e de homologia absolutos e relativos para
os pré-esquemas. Em todos estes exemplos se aplica o corolário da
proposição 3.1, e é precisamente por meio do funtor associado que se deve
definir $\pi_1$ [2]. O mesmo ocorre partindo da categoria dos grupos
algébricos ou quase algébricos sobre um corpo (ou, mais geralmente, sobre um
pré-esquema noetheriano): obtêm-se os «grupos proalgébricos» de Serre [4].

\subsection*{5. Exemplo: «variedades formais»}
\addcontentsline{toc}{subsection}{5. Exemplo: variedades formais}

Sejam $\Lambda$ um anel noetheriano e $\mathcal C$ a categoria das
$\Lambda$-álgebras $A$ que são módulos artinianos de tipo finito sobre
$\Lambda$ (ou, mais brevemente,

\src{375}
as $\Lambda$-álgebras artinianas). São satisfeitas as condições do corolário da
proposição 3.1. A categoria $\operatorname{Pro}(\mathcal C)$ é aqui
equivalente à categoria das álgebras topológicas $O$ sobre $\Lambda$
isomorfas a limites projetivos topológicos
\[
  O=\varprojlim_i O_i
\]
de álgebras $O_i\in\mathcal C$, isto é, cuja topologia é linear, separada e
completa, e tal que para todo ideal aberto $J_i$ de $O$, $O/J_i$ seja uma
álgebra artiniana sobre $\Lambda$. O funtor
$\mathcal C\longrightarrow(\operatorname{Ens})$ associado a tal álgebra não é
outro que
\begin{align*}
  F(A)=h'_O(A)
  &=\{\text{homomorfismos contínuos de $\Lambda$-álgebras topológicas }
       O\longrightarrow A\}\\
  &=\varinjlim_i
    \operatorname{Hom}_{\Lambda\text{-álgebras}}(O_i,A).
\end{align*}

Observe-se, além disso, que a categoria $\mathcal C$ é essencialmente o produto
das categorias análogas correspondentes aos completamentos dos
anéis locais $\Lambda_{\mathfrak m}$, para os ideais maximais
$\mathfrak m$ de $\Lambda$; portanto, se desejado, pode-se restringir ao caso
em que $\Lambda$ seja um anel local completo dessa classe. Em qualquer caso,
$O$ se decompõe canonicamente no produto topológico de componentes
locais, correspondentes aos «pontos» do esquema formal [2] definido por
$O$. Um ponto dessa classe é definido por um elemento $\xi$ de algum
$F(K)$, onde $K\in\mathcal C$ é um corpo (por exemplo, o corpo residual
da componente local considerada); dois pares $(\xi,K)$ e $(\xi',K')$
definem o mesmo ponto se, e somente se, são dominados por um mesmo
$(\xi'',K'')$, ou se dominam um mesmo $(\xi''',K''')$. (Se os
$\Lambda/\mathfrak m$ são algebricamente fechados, basta tomar a soma
disjunta dos $F(\Lambda/\mathfrak m)$.)

É importante dar condições para que a componente local $O_\xi$ de $O$
correspondente a um $\xi\in F(K)$ seja um anel noetheriano. Quando $\Lambda$
é um anel local completo (noetheriano, recordemos), isso equivale a dizer
que $O_\xi$ é isomorfo a um anel quociente de um anel de séries formais
$\Lambda[[t_1,\ldots,t_n]]$ sobre $\Lambda$. Para dar um critério desse tipo,
introduzamos (para todo anel $A$) a $A$-álgebra $I_A$ dos «números
duais» de $A$ mediante
\[
  I_A=A[t]/t^2A[t].
\]
Seja $\varepsilon:I_A\longrightarrow A$ o homomorfismo de aumentação; este define
(se $A\in\mathcal C$) uma aplicação
\[
  F(\varepsilon):F(I_A)\longrightarrow F(A),
\]

\src{376}
e, utilizando que $F$ é exato à esquerda, define-se de maneira
intrínseca uma estrutura de $A$-módulo na parte
\[
  F(I_A,\xi)=F(\varepsilon)^{-1}(\xi)
\]
de $F(I_A)$ formada pelos $\xi'\in F(I_A)$ que «se reduzem a $\xi$»;
utilizando a forma explícita de $F$ em termos de $O$, verifica-se que este
$K$-módulo se identifica com
\[
  \operatorname{Hom}_{\Lambda}(\mathfrak m_\xi/\mathfrak m_\xi^2,A),
\]
onde $\mathfrak m_\xi$ é o núcleo do homomorfismo
$\xi:O\longrightarrow A$, isto é, se $A$ é um corpo, o ideal maximal da
componente local $O_\xi$ de $O$. Deduz-se imediatamente a proposição
seguinte:

\result{PROPOSIÇÃO}{5.1} Seja $\xi\in F(K)$, onde $K\in\mathcal C$ é um
corpo. Para que a componente local correspondente $O_\xi$ de $O$ seja um
anel noetheriano, é necessário e suficiente que o conjunto $F(I_K,\xi)$ de
elementos de $F(I_K)$ que se reduzem a $\xi$ seja um espaço vetorial
de dimensão finita sobre $K$. Nessas condições, tem-se um isomorfismo
canônico
\[
  F(I_K,\xi)=
  \operatorname{Hom}\!\left(
    \frac{\mathfrak m_\xi/\mathfrak m_\xi^2}
         {\operatorname{im}(\mathfrak n_\xi/\mathfrak n_\xi^2)},K
  \right),
\]
onde $\mathfrak n_\xi$ é o ideal maximal de $\Lambda$, núcleo do
homomorfismo $\Lambda\longrightarrow K$.
\footnote{O texto impresso dá aqui uma fórmula que só é correta quando
$\Lambda$ é um corpo. A errata de 1962 prescreve, no caso geral,
o quociente de $\mathfrak m_\xi/\mathfrak m_\xi^2$ pela imagem de
$\mathfrak n_\xi/\mathfrak n_\xi^2$.}
Em particular, a dimensão do $K$-espaço vetorial $F(I_K,\xi)$ é igual à
dimensão do espaço vetorial assim corrigido sobre o corpo
$O_\xi/\mathfrak m_\xi=K(\xi)$.

Suponhamos que $O_\xi$ seja noetheriano e, para simplificar a escrita, que
$\Lambda$ seja local completo e $O=O_\xi$. Diz-se que $O$ é \emph{liso sobre
$\Lambda$} se $O$ é uma álgebra finita e étale sobre a álgebra completada do
localizado de $\Lambda[t_1,\ldots,t_n]$ em um ideal maximal desse anel que
induz o ideal maximal de $\Lambda$; quando a extensão residual de $O$ sobre
$\Lambda$ é trivial (por exemplo, se o corpo residual de $\Lambda$ é
algebricamente fechado), isso equivale a dizer que o próprio $O$ é isomorfo a uma
álgebra de séries formais desse tipo.
\footnote{A errata de 1962 esclarece que esta definição de «liso sobre
$\Lambda$» só é correta quando a extensão residual $k'/k$ é separável;
no caso geral, remete a SGA III, 1.1.}
Finalmente, se já não se supõe necessariamente que $O$ seja noetheriano, dir-se-á
também que $O$ é liso sobre $\Lambda$ se $O$ é isomorfo a um limite
projetivo topológico de $\Lambda$-álgebras quociente que são noetherianas e
lisas sobre $\Lambda$ no sentido anterior. Isso se generaliza imediatamente ao caso
em que $\Lambda$ e $O$ já não são supostos locais. Dito isso, tem-se a
proposição seguinte:

\result{PROPOSIÇÃO}{5.2} Para que $O$ seja liso sobre $\Lambda$, é necessário
e suficiente que o funtor associado $F$ transforme epimorfismos em
epimorfismos.

Isto significa que, para todo homomorfismo sobrejetivo
$A\longrightarrow A'$ em $\mathcal C$, $F(A)\longrightarrow F(A')$ é
sobrejetivo. Naturalmente, basta verificar esta condição quando $A$ é local
e, procedendo passo a passo, quando o ideal de $A$ que é núcleo de
$A\longrightarrow A'$ é anulado pelo ideal maximal de $A$. Na prática,
isso reduz a questão a verificar que certa obstrução, ligada a invariantes
infinitesimais

\src{377}
da situação que dá origem ao funtor $F$, seja nula, problema de natureza
cohomológica.

Para terminar, digamos algumas palavras, no contexto anterior, sobre os
\emph{anéis de definição}. Seja ainda $F$ um funtor de $\mathcal C$ em
$(\operatorname{Ens})$, pro-representável mediante uma $\Lambda$-álgebra
topológica $O$. Então, para todo $A\in\mathcal C$ e todo $\xi\in F(A)$,
existe um menor subanel $A'$ de $A$ tal que $\xi$ seja imagem de um elemento
$\xi'$ de $F(A')$ (que fica então determinado de maneira única): com efeito,
basta interpretar $\xi$ mediante um homomorfismo de $O$ em $A$ e tomar como
$A'$ a imagem de $O$ por este último. Diz-se então que $A'$ é o anel
de definição do objeto $\xi\in F(A)$. Se $u:A\longrightarrow B$ é um
homomorfismo de álgebras e $\eta=F(u)(\xi)$, o anel de definição de $\eta$
é a imagem por $u$ do anel de definição de $\xi$. Quando se parte de um
funtor $F$ de $\mathcal C$ em $(\operatorname{Ens})$, a existência dos
anéis de definição e as propriedades que acabamos de dar são, com pequenas
ressalvas, equivalentes a que $F$ seja pro-representável; isto é, em
geral estão longe de ser triviais.

\fsection{B}{Os dois teoremas de existência}

Conservemos as notações de A, parágrafo 5, e partamos de um funtor covariante
\[
  F:\mathcal C\longrightarrow(\operatorname{Ens});
\]
buscam-se critérios manejáveis para que $F$ seja pro-representável, isto é,
exprimível mediante uma $\Lambda$-álgebra $O$ como acima. Em virtude do
corolário de A, proposição 3.1, para isso é necessário e suficiente que $F$ seja
exato à esquerda. No estado atual da técnica de descida (cf. as
questões formuladas em [3], p. 9), esse critério não pode ser verificado
diretamente sob essa forma nos casos mais importantes, e são necessários
critérios aparentemente menos exigentes.

\result{TEOREMA}{1} Para que o funtor $F$ seja pro-representável, é necessário
e suficiente que satisfaça as duas condições seguintes:
\begin{enumerate}
\item[(i)] $F$ comuta com os produtos finitos;
\item[(ii)] para toda álgebra $A\in\mathcal C$ e todo homomorfismo
$A\longrightarrow A'$ em $\mathcal C$ tal que o diagrama
\[
  A\longrightarrow A'\rightrightarrows A'\otimes_A A'
\]
seja exato ([3], A, definição 1.2), o diagrama transformado
\[
  F(A)\longrightarrow F(A')\rightrightarrows F(A'\otimes_A A')
\]

\src{378}
é exato.
\end{enumerate}
Além disso, basta verificar (ii) quando $A$ é local e, ademais, em um dos
dois casos seguintes:
\begin{enumerate}
\item[(a)] $A'$ é um módulo livre sobre $A$;
\item[(b)] o módulo quociente $A'/A$ é um $A$-módulo de comprimento 1.
\end{enumerate}

A demonstração deste teorema é bastante delicada e não pode ser esboçada aqui.
Limitemo-nos a assinalar que ela repousa essencialmente em um estudo das relações
de equivalência (no sentido das categorias) no espectro de uma álgebra
artiniana (estudo que ainda suscita vários problemas cuja solução parece
indispensável ao desenvolvimento ulterior da teoria).

Nas aplicações, a verificação da condição (i) é sempre trivial.
A verificação de (ii) decompõe-se em duas partes: o caso em que $A'$ é um $A$-módulo livre
corresponde à técnica de descida por morfismos planos ([3], teoremas 1, 2 e
3) e não oferece dificuldade. Para tratar o caso (b), será utilizado o resultado
seguinte:

\result{TEOREMA}{2} Sejam $A$ um anel local artiniano de ideal maximal
$\mathfrak m$ e $A'$ uma $A$-álgebra que contém a $A$, tal que
$\mathfrak mA'\subset A$ e que
\[
  A\longrightarrow A'\rightrightarrows A'\otimes_A A'
\]
seja exato (o que ocorre, em particular, se $A'/A$ é um $A$-módulo de
comprimento 1). Seja $\mathcal F$ a categoria fibrada ([3], A, definição 1.1) de
feixes quase coerentes e planos sobre pré-esquemas variáveis. Então o
morfismo $\operatorname{Spec}(A')\longrightarrow\operatorname{Spec}(A)$ é um
morfismo de $\mathcal F$-descida estrito ([3], A, definição 1.7).

Em outros termos, o dado de um $A$-módulo plano $M$ é completamente
equivalente ao dado de um $A'$-módulo plano $M'$, provido de um
$A'\otimes_AA'$-isomorfismo de $M'\otimes_AA'$ sobre $A'\otimes_AM'$ que
satisfaz a condição usual de transitividade para um dado de descida (loco
citato).

O teorema 2 demonstra-se provando primeiro que
\[
  H^i(A'/A,G_a)=0\qquad\text{para }i\geq 1
\]
([3], A, 4(e)), pois a hipótese $\mathfrak mA'\subset A$ permite reduzir-se
facilmente ao caso em que $A$ é um corpo $k$ (a saber, $A/\mathfrak m$). Se
Aplicam-se então as equivalências assinaladas em [3], página 16.

\fsection{C}{Aplicações a alguns casos particulares}

\subsection*{1. Observações gerais sobre os funtores representados por pré-esquemas}
\addcontentsline{toc}{subsection}{1. Observações gerais sobre os funtores representados por pré-esquemas}

Seja $S$ um pré-esquema localmente noetheriano. Diz-se que um pré-esquema $X$
sobre $S$ é \src{379}localmente de tipo finito sobre $S$ se, para todo
$x\in X$ que se projeta em $y\in S$,\footnote{O texto impresso diz aqui
« $y\in Y$ » e volta a empregar $Y$ ao final do parágrafo. Como o morfismo
considerado é $X\longrightarrow S$, lê-se $S$ em ambas as ocorrências.} existem
uma vizinhança afim de $y$, de anel $A$, e uma vizinhança afim de $x$ situada sobre
o anterior, de anel $B$, tais que $B$ seja uma álgebra de tipo finito sobre
$A$. Encontram-se muitos exemplos importantes de pré-esquemas localmente de
tipo finito sobre $S$ que não são de tipo finito sobre $S$, como soluções de
problemas universais clássicos; assim, importa poder considerar o esquema de
Picard de uma curva como a reunião de uma infinidade de componentes conexas
(que não devem ser confundidas com a componente conexa do elemento neutro, isto
é, a «variedade de Picard»). Em certas ocasiões, é então conveniente
trabalhar na categoria $\mathcal C$ dos pré-esquemas localmente de tipo
finito sobre $S$, para nela examinar a questão da representabilidade de
um funtor contravariante $F$. O objetivo principal destas exposições é
desenvolver uma técnica geral que permita reconhecer se tal funtor $F$ é
representável e estudar as propriedades do $S$-pré-esquema $X$ correspondente
a partir das de $F$. Assinalemos de passagem que, nesse estudo, aparecem
exemplos não patológicos de pré-esquemas sobre $S$ não separados sobre $S$, em
particular como «pré-esquemas de Picard» de excelentes $S$-esquemas; não devem,
portanto, banir-se da geometria algébrica os pré-esquemas que não são
esquemas.

Seja $X$ um pré-esquema localmente de tipo finito sobre $S$, e seja
\[
  F(Y)=\operatorname{Hom}_S(Y,X)
\]
o funtor contravariante associado. Pode-se considerar a restrição $F_0$ de
$F$ à subcategoria $\mathcal C_0$ de $\mathcal C$ formada pelos
pré-esquemas $Y$ sobre $S$ que são artinianos e finitos sobre $S$: quando
$S=\operatorname{Spec}(\Lambda)$, $\mathcal C_0$ é, pois, a categoria dual
da categoria das $\Lambda$-álgebras artinianas considerada em B. Se
$Y=\operatorname{Spec}(A)$, onde $A$ é um anel local artiniano, $Y$
reduz-se a um único ponto $y$ situado sobre um ponto fechado $s$ de $S$, e um
$S$-homomorfismo de $Y$ em $X$ (isto é, um elemento de $F(Y)$) fica definido
pelo dado de um ponto $x\in X$ sobre $s$ e de um
$\mathcal O_{s,S}$-homomorfismo de $\mathcal O_{x,X}$ em $A$. Se existe tal
homomorfismo, então $x$ é necessariamente um ponto fechado de $X$ (pois seu
corpo residual é algébrico sobre o de $s$). Isso mostra que a restrição
$F_0$ de $F$ às «$\Lambda$-álgebras artinianas» é pro-representável, e está
representada pela álgebra topológica sobre $\Lambda$ cujas componentes
locais são os completamentos $\widehat{\mathcal O}_{x,X}$ dos anéis locais
de $X$ nos pontos $x$ de $X$ que são fechados e estão situados sobre pontos
fechados de $S$. Isso mostra que apenas o conhecimento de $F_0$ fornece
informação precisa sobre a estrutura de $X$ (a saber, a estrutura dos
completamentos de seus anéis locais nos pontos indicados). Observe-se que,
mesmo quando $S$

\src{380}
é o espectro de um corpo algebricamente fechado, somente mediante a
consideração sistemática de «variedades» $Y$ tais que $\mathcal O_Y$ possa
ter elementos nilpotentes, em particular trabalhando com os espectros de
anéis artinianos locais, pode-se chegar à «formulação correta» dos
problemas universais clássicos e captar seu aspecto «infinitesimal».

Quando se parte de um funtor $F$ dado de antemão e se quer decidir se ele é
representável, o estudo do funtor $F_0$ (com o auxílio dos teoremas 1 e 2)
fornecerá indicações quase completas: ou, como ocorre com frequência
(verificando simplesmente, por exemplo, a natureza dos conjuntos
$F(I_K,\xi)$ e seu comportamento funtorial, cf. A), constata-se que $F_0$ já não
é pro-representável (o que explica o fracasso das tentativas realizadas
até agora para definir de maneira mais ou menos natural variedades de módulos
para a classificação dos fibrados vetoriais de posto $>1$); ou consegue-se
verificar que $F_0$ é representável, mas os espaços vetoriais
$F(I_K,\xi)$ não são de dimensão finita, caso em que será preciso contentar-se com
a solução «formal»; ou, por fim, $F_0$ é representado por um
produto de anéis locais completos noetherianos, o que fornece indícios muito
sólidos de que o próprio $F$ é representável e, unido a propriedades análogas,
porém de natureza mais global, que pretendemos desenvolver mais adiante, bastará
sem dúvida para implicar que assim seja efetivamente. Finalmente, encontram-se
problemas geométricos interessantes (vejam-se os parágrafos 4 e 5 abaixo) nos
quais se dispõe apenas do funtor $F_0$ (que não provém de nenhum funtor «global»
$F$), e nos quais haverá motivos para se dar por satisfeito se se conseguir associar-lhe
um «esquema formal de módulos».

Para terminar estas generalidades, indiquemos como a teoria de esquemas
explica anomalias aparentes, como a superfície de Igusa $V$, cuja «variedade
de Picard» $P$ se reduz a um ponto e para a qual, no entanto, tem-se
\[
  H^1(V,\mathcal O_V)\ne 0 ;
\]
nesse caso, $P$ é um grupo «puramente infinitesimal» que não se reduz ao
grupo unidade, isto é, está definido por uma álgebra local $O$ de posto finito
sobre o corpo base $k$ e dotado de uma aplicação diagonal correspondente à
estrutura multiplicativa de $P$; se $\mathfrak m$ é o ideal maximal de
$O$, o dual de $\mathfrak m/\mathfrak m^2$ é canonicamente isomorfo a
$H^1(V,\mathcal O_V)$ (cf. o parágrafo 3 mais abaixo). Somente quando o grupo de
Picard é um grupo algébrico no sentido clássico, isto é, liso sobre o
corpo base $k$, a dimensão de $H^1(V,\mathcal O_V)$ (sempre igual à de
$\mathfrak m/\mathfrak m^2$) é igual à do grupo de Picard.

\subsection*{2. Esquemas $\underline{\operatorname{Hom}}_S(X,Y)$,
$\Pi_{X/S}Z$, $\underline{\operatorname{Aut}}_S(X)$, etc.}
\addcontentsline{toc}{subsection}{2. Esquemas Hom, Pi e Aut}

Sejam $X,Y$ dois pré-esquemas sobre $S$; para todo pré-esquema $T$ sobre $S$,

\src{381}
sejam $X_T=X\times_ST$, $Y_T=Y\times_ST$, e tomemos o conjunto
\[
  F(T)=\operatorname{Hom}_T(X_T,Y_T)
      =\operatorname{Hom}_S(X_T,Y)
      =\operatorname{Hom}_S(X\times_ST,Y)
\]
como funtor contravariante em $T$. Se ele for representável, denota-se por
$\underline{\operatorname{Hom}}_S(X,Y)$ o pré-esquema sobre $S$ que o
representa; tem-se, portanto, um isomorfismo funtorial
\[
  \operatorname{Hom}_S
  \bigl(T,\underline{\operatorname{Hom}}_S(X,Y)\bigr)
  \xrightarrow{\sim}
  \operatorname{Hom}_S(T\times_SX,Y).
\]
Há muitas variantes deste problema universal cuja solução se reduz a ele:
o pré-esquema dos $S$-automorfismos de um $S$-pré-esquema $X$ (que será um
pré-esquema em grupos), o pré-esquema dos $S$-homomorfismos de um
$S$-pré-esquema em grupos em outro (que será um pré-esquema em grupos comutativos
se o segundo esquema em grupos é comutativo), etc. Também se pode
generalizar a definição de $\underline{\operatorname{Hom}}_S(X,Y)$
considerando um pré-esquema $Z$ sobre o pré-esquema $X$ sobre $S$, e o funtor
\[
  F(T)=\operatorname{Hom}_{X_T}(X_T,Z_T)
\]
(conjunto das «seções» de $Z_T$ sobre $X_T$); se esse funtor for
representável, o $S$-pré-esquema que o representa será denotado por
$\Pi_{X/S}Z$; ter-se-á, pois, por definição, um isomorfismo funtorial
\[
  \operatorname{Hom}_S(T,\Pi_{X/S}Z)
  =\operatorname{Hom}_{X_T}(X_T,Z_T).
\]
Tomando $Z=Y\times_SX$, recupera-se
$\underline{\operatorname{Hom}}_S(X,Y)$. Dessas definições resulta para os
novos pré-esquemas assim introduzidos um formulário tão trivial quanto útil, que não
daremos aqui (visto que é válido em toda categoria na qual existem
produtos e produtos fibrados). Mais séria é a questão da existência de
esquemas do tipo $\underline{\operatorname{Hom}}_S(X,Y)$. Constata-se primeiro
que, para $X$ fixo, $\underline{\operatorname{Hom}}_S(X,Y)$ somente pode existir
para todo $Y$ sobre $S$ (ou $\Pi_{X/S}Z$ para todo $Z$ sobre $X$) se $X$ é plano
sobre $S$. Além disso, percebe-se que só é razoável esperar a existência de
uma solução, para $Y$ suficientemente geral, se $X$ é além disso próprio sobre
$S$. Por outro lado, parece que essas condições são suficientes para a
existência de $\underline{\operatorname{Hom}}_S(X,Y)$ e $\Pi_{X/S}Z$, com a
condição de impor, quando couber, uma hipótese do tipo «quase-projetivo»
sobre $Y/S$, respectivamente sobre $Z/X$; isso pode ser verificado, em qualquer
caso, em bastantes situações (por exemplo, quando $Y$ é afim sobre $S$; ou,
mediante construções diretas elementares, quando $X$ é finito sobre $S$).
Eis o que fornecem os teoremas 1 e 2:

\src{382}
\result{PROPOSIÇÃO}{2.1} Sejam $\Lambda$ um anel noetheriano, $X$ e $Y$ dois
pré-esquemas quaisquer sobre $\Lambda$, e tomemos o funtor
\[
  F(A)=\operatorname{Hom}_A(X_A,Y_A)
\]
na categoria $\mathcal C_0$ das $\Lambda$-álgebras artinianas. Se $X$ é
plano sobre $\Lambda$, este funtor é pro-representável.

Além disso, verifica-se que para todo $A\in\mathcal C_0$ e $\xi\in F(A)$ tem-se
um isomorfismo canônico
\[
  F(I_A,\xi)=H^1\!\left(
    X_A,
    \underline{\operatorname{Hom}}_{\mathcal O_{X_A}}
    \bigl(\xi^*(\Omega^1_{Y_A/A}),\mathcal O_{X_A}\bigr)
  \right),
\]
onde $\Omega^1_{Y_A/A}$ é o feixe de $1$-diferenciais de Kähler de $Y_A$
com respeito a $A$. Tomando $A$ como um corpo, obtém-se, utilizando A.5.1 e o
teorema de finitude de [2], o corolário seguinte:

\result{COROLÁRIO}{} Suponhamos que $X$ seja plano e próprio sobre $S$, e que $Y$
seja localmente de tipo finito sobre $S$. Então $F$ é pro-representável e as
componentes locais da álgebra topológica correspondente sobre $\Lambda$ são
anéis noetherianos.

\result{OBSERVAÇÕES}{} Os problemas considerados nesta seção, e muitos
outros, eram habitualmente estudados, no âmbito da geometria algébrica
clássica, mediante as «coordenadas de Chow» dos ciclos no espaço
projetivo, que permitiam considerar esses ciclos como pontos de variedades
projetivas adequadas. Esse procedimento, e em geral o uso das
coordenadas de Chow, parece irremediavelmente insuficiente do ponto de
vista dos esquemas, pois destrói os elementos nilpotentes nas
variedades de parâmetros e, em particular, não se presta a um estudo
satisfatório das variações infinitesimais de ciclos (sem falar de sua
natureza não intrínseca, ligada ao espaço projetivo). A linguagem das
coordenadas de Chow era, infelizmente, a única utilizada por muitos
geômetras algébricos para estudar famílias de variedades ou famílias de
ciclos, o que parece ter constituído um sério obstáculo para esclarecer essas
noções, apesar de seu indubitável interesse técnico (provavelmente provisório). Se
se quiser obter o análogo das variedades de Chow na teoria dos
esquemas, chega-se ao seguinte problema universal: seja $X$ um pré-esquema sobre
$S$; para todo pré-esquema $T$ sobre $S$, considera-se o conjunto $F(T)$ dos
subesquemas fechados de $X_T=X\times_ST$ que são planos sobre $T$. Tenta-se
representar este funtor em $T$ mediante um pré-esquema sobre $S$. Mais
geralmente, pode-se dar um feixe quase coerente $G$ sobre $X$ e tomar como
$F(T)$ o conjunto dos feixes

\src{383}
quociente de $G_T$ que são planos sobre $T$. Parece que o problema tem uma
solução, mediante um esquema $C$ localmente de tipo finito sobre $S$, quando
$X$ é próprio sobre o pré-esquema localmente noetheriano $S$ e $G$ é além disso
coerente. Em qualquer caso, supondo unicamente que $S$ seja localmente
noetheriano, a restrição de $F$ às «$S$-álgebras artinianas» é
pro-representável; e se, além disso, $X$ é próprio sobre $S$ e $G$ é coerente, as
componentes locais do anel topológico correspondente $O$ são
noetherianas. Naturalmente, mesmo depois de demonstrar a existência do
«esquema de Chow» de $X$ sobre $S$, restará
encontrar uma decomposição em abertos disjuntos $C_i$ (correspondentes à
fixação de invariantes contínuos, como o grau e a dimensão dos
ciclos que se fazem variar) que sejam de tipo finito (e não somente localmente de
tipo finito) sobre $S$, precisar as relações deste esquema com as
variedades de Chow clássicas e precisar quando um $C_i$ é inclusive projetivo,
ou ao menos quase-projetivo, sobre $S$.\footnote{A errata de 1962
esclarece que o caso projetivo fica completamente resolvido pelos esquemas de
Hilbert (TDTE IV), enquanto exemplos de Nagata e Hironaka mostram que, sem
uma hipótese de projetividade, esses funtores podem não ser representáveis,
mesmo para as subvariedades de dimensão zero de uma variedade completa não
singular de dimensão três; esse fenômeno está ligado à possível
inexistência do quadrado simétrico.}

\subsection*{3. Esquemas de Picard}
\addcontentsline{toc}{subsection}{3. Esquemas de Picard}

Seja $f:X\longrightarrow S$ um $S$-pré-esquema,\footnote{Para um estudo mais
completo, a errata de 1962 remete a TDTE V.} e tomemos o feixe
multiplicativo $\mathcal O_X^*$ das unidades do feixe estrutural de $X$, e
o grupo
\[
  P(X/S)=H^0\bigl(S,R^1f_*(\mathcal O_X^*)\bigr),
\]
chamado grupo de Picard relativo de $X/S$. Um elemento deste grupo fica,
pois, definido dando uma cobertura aberta $(U_i)$ de $S$ e um feixe invertível
$L_i$ sobre cada $f^{-1}(U_i)$, de modo que, para todo par $(i,j)$,
$L_i|f^{-1}(U_i\cap U_j)$ seja isomorfo a $L_j|f^{-1}(U_i\cap U_j)$, ao menos
localmente sobre $U_i\cap U_j$ (isto é, esses dois feixes sejam «equivalentes»
no sentido de [3], B, 4). Se $X/S$ admite uma seção, então $P(X/S)$ nada mais
é que o conjunto das classes de feixes invertíveis sobre $X/S$ módulo
«equivalência» (loco citato). Ponhamos agora, para todo $T$ sobre $S$,
\[
  F(T)=P(X_T/T) ;
\]
obtém-se um funtor covariante em $T$, que poderá ser chamado funtor de Picard de
$X/S$; se esse funtor for representável, o pré-esquema sobre $S$ que o
representa será chamado pré-esquema de Picard de $X/S$ e será denotado por
$\underline P(X/S)$. Ter-se-á, pois, um isomorfismo de funtores
\[
  \operatorname{Hom}_S\bigl(T,\underline P(X/S)\bigr)
  \xrightarrow{\sim}P(X_T/T).
\]
A formação de pré-esquemas de Picard é compatível com a extensão da
base; em particular, os pré-esquemas de Picard das fibras de $X$ sobre $S$
(que são pré-esquemas sobre os corpos residuais $K(s)$ dos $s\in S$) são
as fibras de $\underline P(X/S)$. Naturalmente, como $P(X_T/T)=F(T)$ é um
grupo comutativo, os pré-esquemas de Picard são pré-esquemas em grupos.
Observe-se também que as jacobianas generalizadas de Rosenlicht são precisamente
as componentes conexas da unidade nos esquemas de Picard de curvas
completas que podem ter singularidades, o que deveria tornar evidente a
maioria de suas propriedades (uma vez demonstrada a existência).

\result{OBSERVAÇÃO}{} A definição adotada aqui só é razoável quando
todo ponto de $S$\footnote{O texto impresso diz «$Y$»; aqui não se definiu
nenhuma base $Y$, e a condição se refere à base $S$.} admite uma vizinhança
aberta $U$ sobre a qual $X$ admite uma seção. No caso geral, é preciso
modificar ligeiramente a definição do funtor de Picard para poder obter
também um teorema de existência.

Aqui, as condições plausíveis para a existência de um pré-esquema de Picard
são as seguintes: $X$ é próprio e plano sobre $S$,
$f_*(\mathcal O_X)=\mathcal O_S$, e $X$ admite localmente uma seção sobre
$S$. Essa condição aparece de maneira natural ao aplicar a técnica de
descida, para eliminar os automorfismos de um feixe invertível $L$ sobre $X$
dotando-o de uma seção marcada sobre a seção $s$ ([3], B, 4). Obtém-se,
em particular:

\src{384}
\result{PROPOSIÇÃO}{3.1} Suponhamos que $X$ seja plano sobre
$S=\operatorname{Spec}(\Lambda)$, com $\Lambda$ noetheriano, e que, para todo
$T$ de tipo finito sobre $S$, tem-se
\[
  f_{T*}(\mathcal O_{X_T})=\mathcal O_T
\]
(se $f$ é próprio e separável,\footnote{A errata de 1962 prescreve
explicitamente a hipótese «$f$ próprio e separável».} ou se $S$ é o espectro
de um corpo, resulta de Künneth que a última condição equivale a
$f_*(\mathcal O_X)=\mathcal O_S$). Então o funtor de Picard de $X/S$ na
categoria das $\Lambda$-álgebras artinianas é pro-representável.

Além disso, tem-se aqui
\[
  F(I_A,\xi)=H^1(X_A,\mathcal O_{X_A}),
\]
em particular:

\result{COROLÁRIO}{} Se $X$ é próprio sobre $S$, então as componentes
locais da $\Lambda$-álgebra topológica $O$ correspondente ao funtor de
Picard são noetherianas.

\result{OBSERVAÇÕES}{} Podem-se generalizar as definições e os resultados
desta seção à classificação dos fibrados principais sobre $X$, tendo como
grupo estrutural um esquema em grupos $G$ sobre $S$ que seja afim, plano sobre
$S$ e comutativo. Quando $G$ não é comutativo, e portanto o fibrado em
grupos adjunto de um fibrado principal (cujas seções são os automorfismos
do fibrado principal) deixa de ser trivial, a proposição 3.1 já não é válida
tal como está. Não obstante, pode-se modificar o problema universal de modo que
continue a ter solução (ao menos, por enquanto, em geometria formal). A regra de
ouro que se deve recordar, no contexto desta \src{385}seção e das
seguintes, e sempre que se busquem «esquemas de módulos» para classes de
objetos definidos apenas a menos de isomorfismo, é sempre a seguinte: eliminar
os possíveis automorfismos dos objetos que se querem classificar mediante
a introdução, se necessário, de estruturas auxiliares (pontos ou
elementos de seções marcadas, fixação de formas diferenciais, etc.), que
serão escolhidas suficientemente inócuas para não modificar substancialmente o
problema inicial.

\subsection*{4. Módulos formais de uma variedade}
\addcontentsline{toc}{subsection}{4. Módulos formais de uma variedade}

Seja $\Lambda$ um anel local noetheriano de corpo residual $k$ (com maior
frequência, $\Lambda$ será igual a $k$, ou a um $p$-anel de Cohen), e seja
$X_0$ um pré-esquema sobre $k$. Para toda $\Lambda$-álgebra artiniana local $A$,
considera-se o conjunto $F(A)$ das classes, salvo isomorfismo, de
$A$-pré-esquemas $X$ planos sobre $A$, dotados de um isomorfismo
\[
  (*)\qquad
  X\otimes_A\kappa(A)\xleftarrow{\sim}
  X_0\otimes_k\kappa(A),
\]
onde $\kappa(A)$ é o corpo residual de $A$; naturalmente, os isomorfismos
entre tais $A$-pré-esquemas planos devem respeitar o isomorfismo precedente
incluído na estrutura. Se $A$ é uma $\Lambda$-álgebra artiniana não
necessariamente local, de componentes locais $A_i$, toma-se como $F(A)$ o
produto dos $F(A_i)$. Desse modo, $F$ torna-se um funtor
multiplicativo em $A$, que poderá ser chamado funtor de módulos para $X_0$ (e
$\Lambda$). Se esse funtor for representável, corresponde-lhe uma
$\Lambda$-álgebra topológica $O$ local, de corpo residual $k$, e o espectro
formal de $O$ será chamado esquema formal de módulos para $X_0$ (e $\Lambda$);
cf. [2] para alguns detalhes sobre essa noção.

Aqui, para aplicar a técnica de descida, os automorfismos «finitos» de $X_0$
são inofensivos, pois não influem na existência de automorfismos (no
sentido precisado acima) dos $A$-pré-esquemas $X$; quando $A$ não se
reduz a um corpo, a condição necessária e suficiente para que $X$ não tenha
nenhum $A$-automorfismo não trivial é que se tenha
\[
  H^0(X_0,\Theta_{X_0/k})=0,
\]
onde $\Theta_{X_0/k}$ denota o feixe das $k$-derivações (= feixe tangente)
de $X_0$. Por outro lado, obtém-se facilmente (ao menos se $X_0$ for liso
sobre $k$)
\[
  F(I_A,\xi)=H^1(X_A,\Theta_{X_A/A}).
\]
Conclui-se então, como de costume:

\src{386}
\result{PROPOSIÇÃO}{4.1} Suponhamos que
$H^0(X_0,\Theta_{X_0/k})=0$. Então existe o esquema formal de módulos para
$X_0$. Se, além disso, $X_0$ for próprio sobre $k$, o esquema formal de módulos é
noetheriano.

\result{OBSERVAÇÕES}{} 1\textsuperscript{o} Quando não se supõe que $X_0$ seja
liso sobre $k$, $F(I_A,\xi)$ identifica-se com um sub-$A$-módulo de
\[
  \operatorname{Ext}^1_{\mathcal O_{P_A}}
  \bigl(P_A;\mathcal I_{X_A},\mathcal O_{X_A}\bigr),
\]
onde se põe $P_A=X_A\times_\Lambda X_A$, onde $\mathcal O_{X_A}$ se
considera como um feixe coerente sobre $P_A$ mediante o morfismo diagonal
$X_A\longrightarrow P_A$, e onde $\mathcal I_{X_A}$ denota o feixe coerente
de ideais sobre $P_A$ definido pelo morfismo diagonal. Mais precisamente,
uma globalização simples da teoria de Hochschild mostra que o
$\operatorname{Ext}^1$ anterior identifica-se com o conjunto das classes,
salvo isomorfismo, de feixes de $I_A$-álgebras planas $\mathcal O$ sobre $X_A$,
dotados de um isomorfismo de aumentação
\[
  \mathcal O\otimes_{I_A}A\longrightarrow\mathcal O_{X_A}
\]
(recorde-se que se põe $I_A=A[t]/(t^2)$). O submódulo $F(I_A,\xi)$ é o
que corresponde aos feixes de álgebras comutativas. As hipóteses de lisura
não são, pois, essenciais na teoria de módulos, contrariamente ao que [2]
dava a entender.

2\textsuperscript{o} Recordemos (loco citato) que, em particular, toda curva
algébrica $X_0$ lisa e própria sobre $k$ admite um esquema formal de módulos,
liso sobre $\Lambda$, de dimensão relativa $3g-3$ se o gênero $g$ é
$\geq2$, e $g$ se $g=0$ ou $1$. Esses dois últimos casos já não entram diretamente
na proposição 4.1. No entanto, no caso das curvas elípticas
($g=1$), poderiam ser reduzidos a ela mediante as considerações que seguem.

Naturalmente, pode-se variar ad libitum a proposição 4.1 considerando
sistemas de esquemas sobre $k$ dotados de diversas estruturas. Suponhamos,
por exemplo, que $X_0$ seja um esquema abeliano sobre $k$, com ponto de origem
marcado (isto é, que $X_0$ se considere como esquema em grupos sobre $k$), e
seja $F(A)$ o conjunto das classes, salvo isomorfismo, de esquemas abelianos
sobre $A$ (isto é, de esquemas em grupos próprios e lisos sobre $A$) dotados
de um isomorfismo de esquemas abelianos $(*)$. Verifica-se que a imposição
de uma estrutura multiplicativa (e inclusive de uma mera «seção unidade»)
elimina os automorfismos infinitesimais e que, por conseguinte, existe um
esquema formal de módulos, correspondente a um anel local completo
noetheriano $O$. Por outro lado, pode-se demonstrar que, se $X$ é um esquema
próprio e liso com fibras «absolutamente conexas» sobre um pré-esquema localmente
noetheriano $S$, toda estrutura multiplicativa sobre $X$ que admita uma
seção unidade é necessariamente associativa e comutativa (sempre que o seja,
ao menos, sobre uma fibra e que $S$ seja conexo), e além disso fica determinada de
maneira única pelo conhecimento

\src{387}
da seção unidade. Além disso, suponhamos que $S$ seja o espectro de um anel
local artiniano $A$ de corpo residual $k$, que $X$ seja próprio sobre $A$ e esteja
dotado de uma seção $s$, e que, por fim, $X\otimes_Ak$ esteja dotado de uma
estrutura de esquema abeliano sobre $k$ cujo elemento neutro é o ponto de
$X\otimes_Ak$ correspondente a $s$; um cálculo simples de obstruções,
unido a um raciocínio devido a Serre, permite demonstrar que existe
efetivamente sobre $X$ uma estrutura multiplicativa que admite a seção
$s$ como seção unidade. (A partir daí, utilizando o «teorema de
existência» de [2] para passar ao caso em que $A$ é noetheriano local completo,
e depois a técnica de descida de [3] para o caso geral, pode-se demonstrar
o enunciado análogo para todo $S$ localmente noetheriano conexo.) Isso prova
que o funtor $F(A)$ considerado aqui é isomorfo ao funtor análogo definido ao
princípio desta seção, prescindindo da estrutura multiplicativa sobre
$X_0$. Segue-se, em particular, que, se $\mathfrak m$ é o ideal maximal de
$O$, $\mathfrak m/\mathfrak m^2$ é canonicamente isomorfo ao dual de
$H^1(X_0,\Theta_{X_0/k})$ e tem, portanto, dimensão $n^2$, onde
$n=\dim X_0$. Seria muito interessante determinar se $O$ é liso sobre $\Lambda$,
isto é, isomorfo a uma álgebra de séries formais em $n^2$ indeterminadas
sobre $\Lambda$. O parágrafo 1, proposição 5.2, permite dar uma formulação
equivalente deste problema como problema de existência de esquemas abelianos
que se reduzam a um esquema abeliano dado. Em qualquer caso, por métodos
transcendentes, vê-se que a resposta é afirmativa se $k$ tem característica
zero. Em característica $p\ne0$, basta evidentemente limitar-se ao caso em que
$\Lambda$ é o anel de vetores de Witt construído sobre um corpo
algebricamente fechado $k$. Talvez tenha chegado o momento para o «funtor
de Greenberg» demonstrar seu valor\ldots\footnote{A errata de 1962
anuncia que o esquema formal de módulos de uma variedade abeliana sobre um
corpo é liso sobre o anel de Witt: todo esquema abeliano sobre um quociente
artiniano local se levanta ao anel do qual é quociente. A demonstração
utiliza a variância da classe de obstrução ao levantamento introduzida na
exposição 182, p. 12. Também recorda as construções de Mumford para
as curvas de gênero $g$ e os esquemas abelianos polarizados (cf. SHMT).}

\subsection*{5. Extensão dos recobrimentos}
\addcontentsline{toc}{subsection}{5. Extensão dos recobrimentos}

Sejam $\mathfrak X$ um pré-esquema formal noetheriano [2] e $U$ uma parte
aberta de $\mathfrak X$, definida localmente pela «não anulação» de uma
seção de $\mathcal O_{\mathfrak X}$ que não seja divisor de zero não somente sobre
$\mathfrak X$, mas também sobre todo $X_n$,\footnote{Precisão imposta pela
errata de 1962.} e que seja, portanto, grande o bastante para que toda seção
de $\mathcal O_{\mathfrak X}$ sobre um aberto $V$ que seja nula sobre $U\cap V$
seja nula. Seja $\mathcal J$ um «ideal de definição» para $\mathfrak X$, e seja
\[
  X_n=(\mathfrak X,\mathcal O_{\mathfrak X}/\mathcal J^{n+1}),
\]
que é, portanto, um pré-esquema noetheriano ordinário. Se então
$\mathfrak X',\mathfrak X''$ são dois recobrimentos planos de $\mathfrak X$ (isto
é, dois pré-esquemas sobre $\mathfrak X$ definidos por feixes de álgebras que
são coerentes e localmente livres como feixes de módulos), não ramificados sobre
$U$, a aplicação evidente
\[
  \operatorname{Hom}_{\mathfrak X}(\mathfrak X',\mathfrak X'')
  \longrightarrow
  \operatorname{Hom}_{X_0}(X'_0,X''_0)
\]
é injetiva; em particular, um automorfismo de $\mathfrak X'$ que induz a
identidade sobre $X'_0$ é a identidade. Isso permite aplicar à situação a
técnica de descida. Partamos, em particular, de um recobrimento plano $X'_0$
de $X_0$, não ramificado

\src{388}
sobre $U_0$; seja $G(\mathfrak X)$ o conjunto das classes, salvo isomorfismo
(que induzam a identidade sobre $X'_0$), de recobrimentos planos $\mathfrak X'$
de $\mathfrak X$ que induzem $X'_0$ sobre $X_0$ (e são, portanto,
necessariamente não ramificados sobre $U$). Define-se do mesmo modo $G(V)$ para
toda parte aberta $V$ de $\mathfrak X$, e mais geralmente $G(\mathfrak Y)$
para todo pré-esquema formal $\mathfrak Y$ sobre $\mathfrak X$. Dito isto, os
resultados de [2] e [3] implicam em primeiro lugar os resultados seguintes:

\noindent a. Quando $V$ é um aberto variável de $\mathfrak X$, os $G(V)$
formam um feixe sobre $\mathfrak X$, seja $G_{\mathfrak X}=G$. A restrição desse
feixe a $U$ é o feixe constante cujas fibras se reduzem a um elemento.

Mais geralmente, a determinação das fibras de $G_{\mathfrak X}$ é uma
questão de anéis locais completos; mais precisamente:

\noindent b. Para todo $x\in\mathfrak X$, tem-se
\[
  G_x=G\bigl(\operatorname{Spec}(\mathcal O_{\mathfrak X,x})\bigr)
  \subset
  G\bigl(\operatorname{Spec}(\widehat{\mathcal O}_{\mathfrak X,x})\bigr),
\]
(= classes, salvo isomorfismo, de álgebras finitas e livres $B$ sobre
$\widehat{\mathcal O}_{\mathfrak X,x}$ dotadas de um isomorfismo de
$B\otimes_{\widehat{\mathcal O}_{\mathfrak X,x}}\mathcal O_{X_0,x}$ com
$\widehat{\mathcal O}'_{0x}$, onde $\mathcal O'_0$ é o feixe de álgebras sobre
$X_0$ que define $X'_0$).\footnote{A fórmula impressa identifica diretamente
$G_x$ com o termo completado e omite o acento circunflexo sobre
$\mathcal O'_{0x}$. A inclusão e o completamento são os prescritos pela
errata de 1962.}

\noindent c. Tem-se um isomorfismo canônico
\[
  G_{\mathfrak X}=\varprojlim_nG_{X_n} ;
\]
em outros termos, para todo aberto $V$ de $\mathfrak X$, tem-se
\[
  G(V)=\varprojlim_nG(V_n).
\]

\noindent d. Suponhamos que $\mathfrak X$ seja obtido de um esquema ordinário
próprio $X$ sobre um anel local noetheriano completo $\Lambda$ que possui um
ideal de definição $\mathfrak m$, tomando o completamento $\mathcal J$-ádico de
$\mathcal O_X$, onde $\mathcal J=\mathfrak m\mathcal O_X$. Então
$G(\mathfrak X)$ é canonicamente isomorfo ao conjunto das classes de
recobrimentos planos do esquema ordinário $X$ que «se reduzem a $X'_0$».

De maneira gráfica, pode-se dizer que (a) e (b) estabelecem as relações
fundamentais entre o aspecto local e o aspecto global do problema, (c) dá
as relações entre o aspecto «finito» e o aspecto «infinitesimal» e, por
último, (d) recorda (sob as condições especificadas) a identidade entre o
aspecto «formal» e o aspecto «algébrico».

Suponhamos agora que $\mathfrak X$ esteja definido sobre um anel local
noetheriano completo $\Lambda$, com
$\mathcal J=\mathfrak m\mathcal O_{\mathfrak X}$ (portanto, $X_0$ é um
pré-esquema sobre $\Lambda/\mathfrak m$). Para toda álgebra $A$ finita sobre
$\Lambda$, ponhamos
\[
  F(A)=G(\mathfrak X\times_\Lambda A).
\]
É um funtor covariante em $A$, com valores na categoria dos conjuntos,
e, por (c), esse funtor fica completamente determinado se for conhecido para $A$
artiniana; dizer que esse funtor é pro-representável, isto é, da forma
\[
  F(A)=\operatorname{Hom}_{\Lambda\text{-algèbres top.}}(O,A),
\]

\src{389}
onde $O$ é uma álgebra topológica sobre $\Lambda$ do tipo considerado em A,
parágrafo 5, equivale a dizer que o mesmo ocorre ao se restringir aos $A$
artinianos. A conjunção dos teoremas 1 e 2 implica então efetivamente:

\result{PROPOSIÇÃO}{5.1} O funtor precedente é pro-representável.

Naturalmente, segundo (a), se $U=\mathfrak X$, $G(\mathfrak Y)$ se reduz a um
único elemento para todo $\mathfrak Y$ sobre $\mathfrak X$, e o funtor $F$
tem então escasso interesse (ter-se-á $O=\Lambda$). Parece que, em quase
qualquer outro caso, o anel local topológico $O$ não é noetheriano. Não
obstante, sua existência põe em evidência, de maneira surpreendente, a
natureza «contínua» do conjunto $G(\mathfrak X)$ das soluções (que
corresponde intuitivamente ao fato de que há uma escolha «contínua» no
modo como a ramificação se propaga ao estender $X'_0$). Compare-se esse
resultado com o ponto de vista de J.-P. Serre [4] na teoria do corpo de
classes local, que põe igualmente em relevo o caráter contínuo do grupo de
Galois topológico da extensão abeliana maximal de um corpo local
«geométrico»; ali, o grupo dual (no sentido de Pontryagin) aparece como um
limite indutivo de grupos algébricos (ou ao menos quase algébricos); também
aqui, a classificação das extensões é dada por «variedades» de
dimensão infinita. Por outro lado, no que precede, pode-se tomar como
$\mathfrak X$ o espectro formal de um anel local completo (do qual
$\Lambda$ será, por exemplo, um subanel de Cohen), e cabe esperar que os
resultados desta seção possam ser utilizados no estudo das extensões
de um anel local completo de dimensão $>1$. Tanto no caso local quanto no
global, talvez permitam formular relações precisas entre os fenômenos de
ramificação superior e fenômenos em característica zero (abordáveis por
métodos transcendentes). Em qualquer caso, a análise preliminar da
proposição 5.1 é o que permite estender ao caso «moderadamente ramificado»
os métodos expostos em [2] para estudar o grupo fundamental e resolver por
métodos transcendentes o «problema dos três pontos».

Para terminar, assinalemos que a situação se simplifica quando $X_0$ tem
dimensão 1: então, em virtude de (a) e (b),
\[
  G(\mathfrak X)\simeq
  \prod_iG\bigl(\operatorname{Spec}(\widehat{\mathcal O}_{\mathfrak X,x_i})\bigr),
\]
onde os $x_i$ são os pontos de $X_0-U$: podem-se dar arbitrariamente as
extensões «locais» nos pontos de ramificação. Além disso, quando $X_0$ é
normal, constata-se que o esquema formal de módulos garantido por 5.1 é
liso sobre $\operatorname{Spec}(\Lambda)$.

\src{390}
\begin{center}
  \textbf{BIBLIOGRAFIA}
\end{center}

\begin{description}
  \item[{[1]}] Grothendieck (Alexander). --- \emph{Sur quelques points
    d'algèbre homologique}, Tohoku Math. J., t. 9, 1957, p. 119--221.
  \item[{[2]}] Grothendieck (Alexander). --- \emph{Géométrie formelle et
    géométrie algébrique}, Séminaire Bourbaki, t. 11, 1958/59, fasc. 3,
    n\textsuperscript{o} 182.
  \item[{[3]}] Grothendieck (Alexander). --- \emph{Technique de descente et
    théorèmes d'existence en géométrie algébrique, I : Généralités, Descente
    par morphismes fidèlement plats}, Séminaire Bourbaki, t. 12, 1959/60,
    fasc. 1, n\textsuperscript{o} 190.
  \item[{[4]}] Serre (Jean-Pierre). --- \emph{Corps locaux et isogénies},
    Séminaire Bourbaki, t. 11, 1958/59, fasc. 3, n\textsuperscript{o} 185.
\end{description}
